\chapter{Patrón MVC}

\section{Flujo y valoración}
Tech-Flow aproxima el patrón Modelo–Vista–Controlador sin formalizarlo: no hay capa de servicio explícita ni inyección de controladores, y los controladores JavaFX hablan directamente con CSV. Una lectura honesta:

\begin{table}[H]\centering\caption{Correspondencia con MVC}\begin{tabular}{p{.2\textwidth}p{.62\textwidth}}\toprule
Capa & Presencia en el código\\\midrule
Vista & FXML en \texttt{com.overbyte.vista} + hojas de estilo.\\
Controlador & \texttt{com.overbyte.controllers.*Controller} (JavaFX + \texttt{FXML}).\\
Modelo & Entidades en \texttt{models} y servicios de dominio en \texttt{user} (reglas de negocio: venta, garantías, caja, inventario).\\
Persistencia & \texttt{services.MotorCSV} (acceso a archivos), no es un ORM: es un DAO crudo.\\\bottomrule\end{tabular}\end{table}

Las ventajas: la lógica de negocio no vive en la vista; el dominio se prueba sin levantar la GUI (gracias al directorio temporal de datos); los añadidos recientes (filtro de categorías, precio consistente) se incorporan tocando pocas piezas. Las debilidades: los controladores repiten fórmulas (tema, tablas, guardado .txt) hasta que se extrajeron los ayudantes \texttt{AyudasUI}, \texttt{InformeTexto} y \texttt{AyudasArchivo}; la lectura de CSV desde la GUI omite una capa DAO intermediaria; no existe consolidación de esquemas a nivel de una única fuente formal. En resumen: MVC aproximado con oleadas de refactor trasladadas hacia servicios compartidos; todavía lejos de un MVC estricto con interfaces y escado por capas.
